\RequirePackage[thinlines]{easybmat} %-- muss aufgrund eines Bugs vor etex und tikz geladen werden
\documentclass[a4paper, twoside, headsepline, index=totoc,toc=listof, fontsize=10pt, cleardoublepage=empty, headinclude, DIV=13, BCOR=13mm, titlepage,]{scrartcl}
\usepackage{scrtime} % KOMA, Uhrzeit ermoeglicht
\usepackage{scrpage2} % wie fancyhdr, nur optimiert auf KOMA-Skript, leicht andere Syntax
\usepackage{etoolbox}
\usepackage{letltxmacro}

% Compiler abhaengige IFs
\usepackage{ifxetex,ifluatex}
\newif\ifxetexorluatex
\ifxetex
  \xetexorluatextrue
\else
  \ifluatex
    \xetexorluatextrue
  \else
    \xetexorluatexfalse
  \fi
\fi

%--Farbdefinitionen
%-- muss vor tikz geladen werden
\usepackage[usenames, table, x11names]{xcolor}
\definecolor{dark_gray}{gray}{0.45}
\definecolor{light_gray}{gray}{0.6}

%--Zum Zeichnen
%-- muss vor polyglossia bzw. babel geladen werden
\usepackage{tikz}
\usepackage{tikz-cd}
\usetikzlibrary{external}
\tikzset{>=latex}
\usetikzlibrary{shapes,arrows,intersections}
\usetikzlibrary{calc,3d}
\usetikzlibrary{decorations.pathreplacing,decorations.markings}
\usetikzlibrary{angles}

%-- Konfiguration von tikzexternalize
\tikzexternalize[prefix=tikz/, up to date check=diff]

\ifxetexorluatex
\tikzset{external/system call={xelatex \tikzexternalcheckshellescape %-- verwende LuaLaTeX, wegen dynamischer Speicherallokation
    -halt-on-error -interaction=batchmode -jobname "\image" "\texsource"}}
\else
\tikzset{external/system call={pdflatex \tikzexternalcheckshellescape 
    -halt-on-error -interaction=batchmode -jobname "\image" "\texsource"}}
\fi

%-- tikzexternalize fuer tikzcd deaktivieren, da inkompatibel
\AtBeginEnvironment{tikzcd}{\tikzexternaldisable}
\AtEndEnvironment{tikzcd}{\tikzexternalenable}

%-- um Inkompatibilitaeten von quotes und polyglossia bzw. babel zu vermeiden
\tikzset{
  every picture/.append style={
    execute at begin picture={\shorthandoff{"}},
    execute at end picture={\shorthandon{"}}
  }
}
\usetikzlibrary{quotes}
\usepackage{pgfplots}
\usepgfplotslibrary{colormaps}
\pgfplotsset{compat=1.10}



%-- Mathe
\usepackage{mathtools} % beinhaltet amsmath
\mathtoolsset{showonlyrefs, centercolon}
\newtagform{brackets}[\textbf]{[}{]}
\usetagform{brackets}
\usepackage{wasysym}
\usepackage{amssymb} %zusätzliche Symbole
\usepackage{latexsym} %zusätzliche Symbole
\usepackage{stmaryrd} %für Blitz
\usepackage{nicefrac} %schräge Brüche
\usepackage{cancel} %Befehle zum Durchstreichen
\usepackage{mathdots}
% \usepackage{bm}
\DeclareSymbolFont{bbold}{U}{bbold}{m}{n}
\DeclareSymbolFontAlphabet{\mathbbold}{bbold}
\newcommand{\mathds}[1]{\mathbb{#1}} %Um Kompatibilitaet mit frueheren Benutzung von dsfont herzustellen
\newcommand{\ind}{\mathbbold{1}} % charakteristische-Funktion-Eins
\def\mathul#1#2{\color{#1}\underline{{\color{black}#2}}\color{black}} %farbiges Untersteichen im Mathe-Modus

%-- Alles was mit Schrift und XeteX zu tun hat
\ifxetexorluatex
    % XeLaTeX or LuaTeX
	\ifxetex
		\usepackage{mathspec} %beinhaltet fontspec 
	\else
		\usepackage[no-math]{fontspec}
	\fi
    \usepackage{polyglossia} %babel-ersatz
    \setmainlanguage[spelling=new,babelshorthands=false]{german}
	\newcommand\glqq{"}
	\newcommand\grqq{"}
	\defaultfontfeatures{Mapping=tex-text, WordSpace={1.4}} %
    \setmainfont[Ligatures=Common, BoldFont={* Bold}, ItalicFont={* Light Italic},SmallCapsFont={Linux Libertine O}, SmallCapsFeatures={Letters=SmallCaps}]{Source Sans Pro}
	\setsansfont[Scale=MatchLowercase,Ligatures=Common, BoldFont={* Medium}]{Ubuntu}
	\ifxetex
		\setallmonofonts[Scale=MatchLowercase, ItalicFont={*}]{Consolas} 
	\else
		\setmonofont[Scale=MatchLowercase, ItalicFont={*}]{Consolas}
	\fi
	\usepackage{xltxtra}
	\usepackage{fontawesome}
	\usepackage{microtype}
\else
    % default: pdfLaTeX
    \usepackage[ngerman]{babel}
    \usepackage[T1]{fontenc}
    \usepackage[utf8]{inputenc}
	\usepackage{textcomp} %verhindert ein paar Fehler bei den Fonts
    \usepackage[babel=true, tracking=true]{microtype}
	\usepackage{lmodern}
	\renewcommand{\familydefault}{\sfdefault} 
\fi


%--Mixed
\usepackage[neverdecrease]{paralist}
\usepackage[german=quotes]{csquotes}
\usepackage{makeidx}
\usepackage{booktabs}
\usepackage{wrapfig}
\usepackage{float}
\usepackage[margin=10pt, font=small, labelfont={sf, bf}, format=plain, indention=1em]{caption}
\captionsetup[wrapfigure]{name=Abb. }
\usepackage{stackrel}
\usepackage{ifthen}
\usepackage{multicol}
\usepackage{qrcode}
\flushbottom

%--Hyphenation
\hyphenation{di-men-si-o-na-le}


%--Unterstreichung
\usepackage[normalem]{ulem}
\setlength{\ULdepth}{1.8pt}

%--Indexverarbeitung
\newcommand{\bet}[1]{\textbf{#1}}
\newcommand{\Index}[1]{\textbf{#1}\index{#1}}
\makeindex
\setindexpreamble{{\noindent \itshape Die \emph{Seitenzahlen} sind mit Hyperlinks zu den entsprechenden Seiten versehen, also anklickbar} \ifxetexorluatex \faHandUp\fi \par \bigskip}
\renewcommand{\indexpagestyle}{scrheadings}


%--Marginnote & todonotes
\usepackage{marginnote}
\renewcommand*{\marginfont}{\color{dark_gray} \itshape \footnotesize }
\usepackage[textsize=small]{todonotes}

%--Todonotes mit tikz/externalize kompatibel machen
\LetLtxMacro{\oldtodo}{\todo}
\renewcommand{\todo}[2][]{\tikzexternaldisable\oldtodo[#1]{#2}\tikzexternalenable}
\LetLtxMacro{\oldmissingfigure}{\missingfigure}
\renewcommand{\missingfigure}[2][]{\tikzexternaldisable\oldmissingfigure[{#1}]{#2}\tikzexternalenable}

%--Konfiguration von Hyperref pdfstartview=FitH, 
\usepackage[hidelinks, pdfpagelabels,  bookmarksopen=true, bookmarksnumbered=true, linkcolor=black, urlcolor=SkyBlue2, plainpages=false,pagebackref, citecolor=black, hypertexnames=true, pdfauthor={Jannes Bantje}, pdfborderstyle={/S/U}, linkbordercolor=SkyBlue2, colorlinks=false]{hyperref}

% \let\orighref\href
\ifxetexorluatex
\newcommand{\hrefsym}[2]{\href{#1}{\texttt{#2\,\faExternalLink}}}
\newcommand{\hrefsymmail}[2]{\href{#1}{\texttt{\faEnvelope\,#2}}}
\renewcommand{\url}[1]{\hrefsym{#1}{\nolinkurl{#1}}}
\fi
\providecommand{\hrefsym}{\href}
% \usepackage{cleveref}
% \crefformat{footnote}{#2\footnotemark[#1]#3}

%--Römische Zahlen
\newcommand{\RM}[1]{\MakeUppercase{\romannumeral #1{}}}

%--Punkte (nach hyperref)
\usepackage{ellipsis}

%%--Abkürzungen etc.
\newcommand{\light}{\text{\Large $\lightning$}}

%-- Definitionen von weiteren Mathe-Befehlen
\DeclareMathOperator{\re}{Re} %Realteil
\DeclareMathOperator{\im}{Im} %Imaginaerteil
\DeclareMathOperator{\id}{id} %identische Abbildung
\DeclareMathOperator{\Sp}{Sp} %Spur
\DeclareMathOperator{\sgn}{sgn} %Signum
\DeclareMathOperator{\alt}{Alt} %Alternierende n-linearFormen
\DeclareMathOperator{\End}{End} %Endomorphismen
\DeclareMathOperator{\Vol}{Vol} %Volumen
\DeclareMathOperator{\dom}{dom} %Domain
\DeclareMathOperator{\Hom}{Hom} %Homomorphismen
\DeclareMathOperator{\bild}{Bild} %Bild
\DeclareMathOperator{\Kern}{Kern }
\DeclareMathOperator{\rg}{rg} %Rang
\DeclareMathOperator{\Rg}{Rg} %Rang
\DeclareMathOperator{\diam}{diam} %Durchmesser
\DeclareMathOperator{\dist}{dist} %Distanz
\DeclareMathOperator{\grad}{grad} %Gradient
\DeclareMathOperator{\dive}{div} %Gradient
\DeclareMathOperator{\rot}{rot} %Rotation
\DeclareMathOperator{\hess}{Hess} %Hesse-Matrix
\DeclareMathOperator{\koker}{Koker} %Kokern
\DeclareMathOperator{\aut}{Aut} %Automorphismen
\DeclareMathOperator{\ord}{ord} %Ordnung
\DeclareMathOperator{\ggT}{ggT} %ggT
\DeclareMathOperator{\kgV}{kgV} %kgV
\DeclareMathOperator{\Gr}{Gr} %Gerade
\DeclareMathOperator{\Kr}{Kr} %Kreis
\DeclareMathOperator{\Char}{char} %Charakteristik
\DeclareMathOperator{\Aut}{Aut} %Automorphismen
% \DeclareMathOperator{\D}{D} %formale Ableitung
\newcommand{\D}{\ensuremath{\mathrm{D}\mkern-1.5mu}}
\newcommand{\mathd}{\ensuremath{\mathrm{d}\mkern-1.0mu}}
\newcommand{\Tmap}{\ensuremath{\mathrm{T}\mkern-1.0mu}}
\newcommand{\opL}{\ensuremath{\mathrm{L}\mkern-0.6mu}}
\DeclareMathOperator{\cov}{cov} %Kovarianz
\DeclareMathOperator{\Gal}{Gal} %Galois
\DeclareMathOperator{\supp}{supp} %Träger
\DeclareMathOperator{\Sym}{Sym} %Symmetrische Gruppe
\DeclareMathOperator{\Zyl}{Zyl} %Zylinder
\DeclareMathOperator{\Mod}{Mod} %Moduln
\DeclareMathOperator{\EPK}{EPK} %Einpunktkompaktifizierung
\DeclareMathOperator{\conj}{conj} %Konjugation
\DeclareMathOperator{\ann}{ann} %Annulator
\DeclareMathOperator{\tor}{tor} %Torsionsmodul
\DeclareMathOperator{\Int}{Int} %Interior/Inneres
\DeclareMathOperator{\Br}{Br} %Brauerergruppe
\DeclareMathOperator{\GL}{GL} %allgemeine lineare Gruppe
\DeclareMathOperator{\SL}{SL} %spezielle lineare Gruppe
\DeclareMathOperator{\SU}{SU} %spezielle unitäre Gruppe
\DeclareMathOperator{\Tr}{Tr} %Spur/Trace
\DeclareMathOperator{\conv}{conv} %Konvexe Teilmenge
% \DeclareMathOperator{\tr}{tr}
\newcommand{\tr}{\mathrm{tr}}
\newcommand{\w}{\mathrm{w}}
\newcommand{\weakT}[1]{\ensuremath{\mathcal{T}_{#1}^{\mkern+1.0mu\text{\raisebox{0.4ex}{$\mathrm{w}$}}}}}
\newcommand{\weakTstar}[1]{\ensuremath{\mathcal{T}_{#1}^{\mkern+1.0mu\text{\raisebox{0.4ex}{$\mathrm{w}$}}^*}}}
\newcommand{\rand}[1]{\ensuremath{\partial^{\scriptscriptstyle #1}}}


% -- Kategorien
\DeclareMathOperator{\Ob}{Ob} %Objekte
\DeclareMathOperator{\Mor}{Mor} %Morphismen

% \DeclareMathOperator{\TOP}{TOP} %Topologische Räume
\newcommand{\TOP}{\textsc{Top}}
% \DeclareMathOperator{\HTOP}{HTOP} %Topologische Räume + Homotopieäquivalenz
\newcommand{\HTOP}{\textsc{HTop}}
% \DeclareMathOperator{\VR}{VR} %Vektorräume
\newcommand{\VR}{\textsc{VR}}
% \DeclareMathOperator{\MOD}{MOD} %Moduln
\newcommand{\MOD}{\textsc{Mod}}
% \DeclareMathOperator{\MENGEN}{MENGEN} %Mengen
\newcommand{\MENGEN}{\textsc{Mengen}}
% \DeclareMathOperator{\MAN}{MAN} %Mannigfaltigkeiten
\newcommand{\MAN}{\textsc{Man}}
% \DeclareMathOperator{\GRUPPEN}{GRUPPEN} %Gruppen
\newcommand{\GRUPPEN}{\textsc{Gruppen}}
% \DeclareMathOperator{\ABELGRUPPEN}{ABEL.GRUPPEN} %abelsche Gruppen
\newcommand{\ABELGRUPPEN}{\textsc{Abel.Gruppen}}
% \DeclareMathOperator{\KAT}{KAT} %Kategorien
\newcommand{\KAT}{\textsc{Kat}}
% \DeclareMathOperator{\FUN}{FUN} %Funktoren
\newcommand{\FUN}{\textsc{Fun}}


\newcommand{\op}{\mathrm{op}}

\newcommand{\Ker}{\ensuremath{\Kern \,}}
\newcommand{\zirkel}{\rotatebox[origin = c]{-90}{$\varangle$}}
\newcommand{\wirkung}{\!\curvearrowright\!}

\newcommand{\oE}{\OE}

%--Beweisende
\newcommand{\bewende}{\ifmmode \tag*{$\square$} \else \hfill $\square$ \fi}

%--Volumen
\newcommand{\vol}[1]{\ensuremath{\Vol_{#1}}}

%--Skalarprodukt
\DeclarePairedDelimiterX\sprod[2]{\langle}{\rangle}{#1\,\delimsize\vert\,#2}

%--Betrag, Gaußklammer
\DeclarePairedDelimiter{\abs}{\lvert}{\rvert}
\DeclarePairedDelimiter{\ceil}{\lceil}{\rceil}

%--Norm
\DeclarePairedDelimiter\doppelstrich{\Vert}{\Vert}
\newcommand{\norm}[2][\relax]{
\ifx#1\relax \ensuremath{\doppelstrich*{#2}}
\else \ensuremath{\doppelstrich*{#2}_{#1}}
\fi}


%--Umklammern
\DeclarePairedDelimiter\enbrace{(}{)}
\DeclarePairedDelimiter\benbrace{[}{]}
\DeclarePairedDelimiter\lenbrace{<}{>}

%--Mengen
\DeclarePairedDelimiterX\mengenA[1]{\lbrace}{\rbrace}{#1}
\DeclarePairedDelimiterX\mengenB[2]{\lbrace}{\rbrace}{#1\, \delimsize\vert \, #2}

\makeatletter
\newcommand{\set}[2][\relax]{
\ifx#1\relax \ensuremath{
\mengenA*{#2}}
\else \ensuremath{%
  \mengenB*{#1}{#2}}
\fi}
\makeatother


%--offen und abgeschlossen
\newcommand{\off}{\! \stackrel[\text{offen}]{}{\subset} \!}
\newcommand{\abg}{\! \stackrel[\text{abg.}]{}{\subset} \!}

%--Differential
\newcommand{\diff}[2]{\ensuremath{\frac{{\partial #1}}{{\partial #2}} }}
\newcommand{\diffd}[2]{\ensuremath{\frac{\mathd #1}{\mathd #2} }}

%--Diagonale Linien für Matrizen
\newcommand\tikzmark[1]{%
  \tikzexternaldisable\tikz[overlay,remember picture,baseline] \node [anchor=base] (#1) {};\tikzexternalenable}

\newcommand\MyLine[3][]{%
  \tikzexternaldisable\begin{tikzpicture}[overlay,remember picture]
    \draw[#1] (#2.south) -- (#3.north);
  \end{tikzpicture}\tikzexternalenable}
  
%--Jordankasten
\newcommand{\Jord}{ \tikzexternaldisable \begin{BMAT}{|ccc|}{|ccc|} %
				0 & & \\ %
				1\tikzmark{z} &  \diagdown & \\ %
				& \tikzmark{u}1 \MyLine[thick]{z}{u} &  0 %
			\end{BMAT} \tikzexternalenable}
			
%--Jordankasten mit Lambda
\newcommand{\JordLa}[1][\empty]{%
	\tikzexternaldisable\ifthenelse{\equal{#1}{\empty}}
		{\begin{BMAT}{|ccc|}{|ccc|} %
				\lambda & & \\ %
				1\tikzmark{z} &  \diagdown & \\ %
				& \tikzmark{u}1 \MyLine[thick]{z}{u} &  \lambda  %
			\end{BMAT}}
		{\begin{BMAT}{|ccc|}{|ccc|} %
				\lambda_{#1}  & & \\ %
				1\tikzmark{z} &  \diagdown & \\ %
				& \tikzmark{u}1 \MyLine[thick]{z}{u} &  \lambda_{#1}  %
			\end{BMAT}}\tikzexternalenable}

%--Inhaltsverzeichnis
\usepackage[tocindentauto]{tocstyle}
\usetocstyle{KOMAlike}			

\newcommand{\fach}{Topologie \RM{1}.}
% \newcommand{\shortFach}{Analysis, Topologie, Geometrie}
\newcommand{\semester}{WiSe 2014/2015}
\newcommand{\homepage}{https://wwwmath.uni-muenster.de/reine/u/topos/lehre/WS2014-2015/Topologie1/Topologie1.html}

\newcommand{\prof}{Prof.\,Dr.\,Arthur Bartels}
\providecommand{\verfasser}{Jannes Bantje}

%--Konfiguration von scrheadings
\setheadsepline{1pt}[\color{light_gray}]
\pagestyle{scrheadings}
%\fancyhf[H, F]{}
\clearscrheadfoot

\providecommand{\shortFach}{\fach}
\lehead{ \includegraphics[height=0.5 cm,keepaspectratio]{../!config/Bilder/Logo_WWU_Muenster_light_gray.pdf}}
\rehead{\rule{0cm}{0.5cm}\footnotesize \sffamily \color{light_gray} \verfasser -- Mitschrift \shortFach}
\lohead{\rule{0cm}{0.5cm} \footnotesize \sffamily \color{light_gray} Stand: \today \; \thistime[:]}
\rohead{\includegraphics[height=0.5 cm,keepaspectratio]{../!config/Bilder/fb10logo_gray.pdf}}


\ofoot[{ \color{dark_gray} \LARGE \sffamily \thepage}]{{ \color{dark_gray} \LARGE \sffamily \thepage}} %hier wir auch der plain Stil bearbeitet!
\automark{section}
\ifoot{ \color{dark_gray} \small \leftmark}

%--Metadaten
\author{\verfasser}
\titlehead{\includegraphics[height=1.5cm, keepaspectratio]{../!config/Bilder/Logo_WWU_Muenster.pdf}%
\hfill \includegraphics[height=1.3cm, keepaspectratio]{../!config/Bilder/fb10logo.pdf}}
\title{Skript \fach}
\subtitle{Mitschrift der Vorlesung  \enquote{\fach} von \prof}

\publishers{
Aktuelle Version verfügbar bei: \bigskip\\
\normalsize
\begin{minipage}[t]{0.45\textwidth}
	\centering{
	\qrcode[height=3.5cm, version=4]{https://github.com/JaMeZ-B/latex-wwu}
	\medskip\\
	\includegraphics[height=0.4cm, keepaspectratio]{../!config/Bilder/github_octo.pdf} 
	\includegraphics[height=0.4cm, keepaspectratio]{../!config/Bilder/GitHub_Logo.pdf} (inklusive Sourcecode)\\
	{\small\url{https://github.com/JaMeZ-B/latex-wwu}}}
\end{minipage}
\hfill
\begin{minipage}[t]{0.45\textwidth}
	\centering{
	\qrcode[height=3.5cm, version=4]{btsync://B6WH2DISQ5QVYIRYIEZSF4ZR2IDVKPN3I?n=latex_share}
	\medskip\\
	\includegraphics[height=0.4cm, keepaspectratio]{../!config/Bilder/bt_sync_logo.pdf} 
	{\large \textbf{Bittorrent} Sync} \\
	\texttt{B6WH2DISQ5QVYIRYIEZSF4ZR2IDVKPN3I}
	}
\end{minipage}
}


\begin{document}
\pagenumbering{Roman}
\maketitle
\begin{abstract}
\section*{Vorwort --- Mitarbeit am Skript}
Dieses Dokument ist eine Mitschrift aus der Vorlesung \enquote{\fach, \semester}, gelesen von \prof. Der Inhalt entspricht weitestgehend dem Tafelanschrieb. Für die
Korrektheit des Inhalts übernehme ich keinerlei Garantie! Für Bemerkungen und Korrekturen -- und seien es nur Rechtschreibfehler -- bin ich sehr dankbar. 
Korrekturen lassen sich prinzipiell auf drei Wegen einreichen: 
\begin{itemize}
	\item Persönliches Ansprechen in der Uni, Mails an \hrefsymmail{mailto:\mail}{\mail} (gerne auch mit annotieren PDFs) 
	\item \emph{Direktes} Mitarbeiten am Skript: Den Quellcode poste ich auf GitHub (siehe oben), also stehen vielfältige Möglichkeiten der Zusammenarbeit zur Verfügung:
	Zum Beispiel durch Kommentare am Code über die Website und die Kombination Fork + Pull Request. Wer sich verdient macht oder ein Skript zu einer Vorlesung, die 
	ich nicht besuche, beisteuern will, dem gewähre ich gerne auch Schreibzugriff.
	
	Beachten sollte man dabei, dass dazu ein Account bei \url{github.com} notwendig ist, der allerdings ohne Angabe von persönlichen Daten angelegt werden kann. 
	Wer bei GitHub (bzw. dem zugrunde liegenden Open-Source-Programm \enquote{\texttt{git}}) -- verständlicherweise -- Hilfe beim Einstieg braucht, dem helfe ich gerne 
	weiter. Es gibt aber auch zahlreiche empfehlenswerte Tutorials im Internet.\footnote{zB. \url{https://try.github.io/levels/1/challenges/1}, ist auf Englisch, aber dafür 
	interaktives LearningByDoing}
	\item \emph{Indirektes} Mitarbeiten: \TeX-Dateien per Mail verschicken. 
	
	Dies ist nur dann sinnvoll, wenn man einen ganzen Abschnitt ändern möchte (zB. einen alternativen Beweis geben), da ich die Änderungen dann per Hand einbauen muss! Ich freue mich aber auch über solche Beiträge!
\end{itemize}
\section*{Vorlesungshomepage}
{\centering 
\begin{minipage}[c][][c]{\textwidth}
	\centering \qrcode[height=3cm]{\homepage} \medskip\\
	\footnotesize \url{\homepage}
\end{minipage}
\par}

\end{abstract}

\setcounter{page}{1}
\tableofcontents
\cleardoubleoddemptypage

\pagenumbering{arabic}
\setcounter{page}{1}

\section{Kategorien, Funktoren und natürliche Transformationen} % (fold)
\label{sec:1}

\subsection[Definition: Kategorie]{Definition} % (fold)
\label{sub:11}
Eine \Index{Kategorie} $\mathcal{C}$ besteht aus: \marginnote{13. Okt}
\begin{enumerate}[i)]
	\item Einer Klasse $\Ob(\mathcal{C})$. Die Elemente von $\Ob(\mathcal{C})$ heißen die \Index{Objekte} von $\mathcal{C}$.
	\item Zu je zwei Objekten $A,B$ aus $\mathcal{C}$ einer Menge $\Mor_{\mathcal{C}}(A,B)$. Die Elemente von $\Mor_\mathcal{C}(A,B)$ heißen \Index{Morphismen} von
	$\mathcal{C}$.
	\item Zu je drei Objekten $A,B,C$ aus $\mathcal{C}$ einer Abbildung: 
	\[
		\Mor_\mathcal{C}(B,C) \times \Mor_\mathcal{C}(A,B) \to \Mor_\mathcal{C}(A,C) \enspace, \quad (f,g) \mapsto f \circ g
	\]
	genannt die \Index{Komposition} in $\mathcal{C}$. 
\end{enumerate}
Dabei müssen folgende Axiome  erfüllt sein:
\begin{enumerate}[(i)]
	\item Die Komposition ist \Index{assoziativ}: \\ Für Objekte $A,B,C,D$ von $\mathcal{C}$ und 
	$f \in \Mor_\mathcal{C}(C,D), g \in \Mor_\mathcal{C}(B,C), h \in \Mor_\mathcal{C}(A,B)$ gilt immer
	\[
		f \circ (g \circ h) = (f \circ g) \circ h
	\]
	\item Die Komposition ist \Index{unital}: Für jedes Objekt $A$ von $\mathcal{C}$ gibt es einen Morphismus $\id_A \in \Mor_\mathcal{C}(A,A)$ so, dass
	\[
		\forall f \in \Mor_\mathcal{C}(B,A) : \id_A \circ f = f  \qquad  \text{und} \qquad  \forall f \in \Mor_\mathcal{C}(A,B) : f \circ \id_A = f
	\]
\end{enumerate}
% subsection 11 (end)

\subsection[Beispiele für Kategorien]{Beispiele} % (fold)
\label{sub:12}
\begin{enumerate}[(1)]
	\item Die Kategorie der Mengen $\MENGEN$: Objekte sind Mengen und Morphismen sind Abbildungen.
	\item Die Kategorie der Gruppen $\GRUPPEN$: Objekte sind Gruppen und Morphismen sind Gruppenhomomorphismen.
	\item Die Kategorie der $K$-Vektorräume $K\text{-}\VR$: Objekte sind $K$-Vektorräume und die Morphismen sind $K$-lineare Abbildungen.
	\item Die Kategorie der $R$-Moduln $R\text{-}\MOD$: Objekte sind $R$-Moduln und Morphismen sind $R$-lineare Abbildungen.
	\item Die Kategorie der $C^\infty$-Mannigfaltigkeiten $C^\infty\text{-}\MAN$: Objekte sind $C^\infty$-Mannigfaltigkeiten und Morphismen sind $C^\infty$-Abbildungen.
	\item Die Kategorie der topologischen Räume $\TOP$: Objekte sind topologische Räume und Morphismen sind stetige Abbildungen.
	\item Die Kategorie der punktierten topologischen Räume $\TOP_\bullet$: Objekte sind punktiert-topologische Räume und Morphismen sind punktiert-stetige Abbildungen.
	\item Die Kategorie $\HTOP$: Objekte sind topologische Räume und Morphismen sind Homotopieklassen von stetigen Abbildungen.
	\item Die Kategorie $\HTOP_\bullet$: Objekte sind punktiert-topologische Räume und Morphismen sind Homotopieklassen von punktiert-stetigen Abbildungen.
	\item Sei $G$ eine Gruppe. Wir erhalten eine Kategorie $\mathcal{C}_G$ mit genau einem Objekt $*$ und $\Mor_{\mathcal{C}_G}(*,*) = G$. Die Komposition wird durch die
	Verknüpfung in der Gruppe festgelegt.
	\item Ist $\mathcal{C}$ eine Kategorie, so ist $\mathcal{C}^{\op}$ eine Kategorie, wobei $\Ob(\mathcal{C}^\op) = \Ob(\mathcal{C})$ und 
	\(
		\Mor_{\mathcal{C}^\op}(A,B) := \Mor_\mathcal{C}(B,A)
	\).
	Die Komposition ist gegeben durch $f \circ^\op g := g \circ f$.
\end{enumerate}
% subsection 12 (end)

\subsection[Bemerkungen zu Kategorien]{Bemerkung} % (fold)
\label{sub:13}
\begin{enumerate}[(i)]
	\item Eine Kategorie heißt \bet{klein}\index{Kategorie!klein}, wenn ihre Objekte eine Menge bilden.
	\item Statt $f \in \Mor_\mathcal{C}(A,B)$ schreiben wir oft $f : A \to B$ oder $A \xrightarrow{f} B$.
\end{enumerate}
% subsection 13 (end)

\subsection[Definition: Funktor]{Definition} % (fold)
\label{sub:14}
Seien $\mathcal{C}$ und $\mathcal{D}$ Kategorien: Ein \Index{Funktor} $F : \mathcal{C} \to \mathcal{D}$ ordnet jedem Objekt $C$ von $\mathcal{C}$ ein Objekt $F(C)$ von 
$\mathcal{D}$ und ordnet jedem Morphismus $f : C \to C'$ in $\mathcal{C}$ einen Morphismus $F(f) : F(C) \to F(C')$ in $\mathcal{D}$ zu. Dabei muss gelten:
\[
	F(f \circ g) = F(f) \circ F(g) \quad \text{und} \quad F(\id_C) = \id_{F(C)}
\] 
% subsection 14 (end)

\subsection[Beispiele für Funktoren]{Beispiele} % (fold)
\label{sub:15}
\begin{enumerate}[(1)]
	\item Es gibt offensichtliche \enquote{Vergiss}-Funktoren:
	\begin{align*}
		K\text{-}\VR \enspace\longrightarrow & \enspace\MENGEN \\
		R\text{-}\MOD \enspace\longrightarrow & \enspace\MENGEN \\
		\TOP \enspace\longrightarrow & \enspace\MENGEN \\
		\TOP_\bullet \enspace\longrightarrow & \enspace\MENGEN \\
		\TOP_\bullet \enspace\longrightarrow & \enspace \TOP
	\end{align*}
	\item Sei $\mathcal{C}$ eine Kategorie und $C$ ein Objekt von $\mathcal{C}$. Der durch $C$ \bet{dargestellte Funktor}\index{Funktor!dargestellter} 
	$F_C : \mathcal{C} \to \MENGEN$ ist definiert durch 
	\[
		\begin{array}{ll}
			F_C(A) = \Mor_\mathcal{C}(C,A) &\text{ für } A \in  \Ob(\mathcal{C}) \\
			F_C(f)  : \Mor_\mathcal{C}(C,A) \to \Mor_\mathcal{C}(C,A'), g \mapsto f \circ  g &\text{ für } f : A \to A' \text{ in } \mathcal{C}
		\end{array}
	\]
	\item Die Fundamentalgruppe definiert einen Funktor
	\[
		\pi_1 : \HTOP_\bullet \longrightarrow \GRUPPEN
	\]
\end{enumerate}
\minisec{Bemerkung}
\begin{itemize}
	\item Ist $F : \mathcal{C} \to \mathcal{D}$ ein Funktor so schreiben wir oft kürzer und ungenauer $f_* := F(f)$.
	\item Kleine Kategorien und Funktoren bilden die Kategorie $\KAT$.
\end{itemize}
% subsection 15 (end)

\subsection[Definition: Natürliche Transformation]{Definition} % (fold)
\label{sub:16}
Seien $F,G : \mathcal{C} \to \mathcal{D}$ zwei Funktoren. Eine \Index{natürliche Transformation} $\tau : F \to G$ ordnet jedem $C \in \Ob(\mathcal{C})$ einen Morphismus
$\tau_C : F(C) \to G(C)$ in $\mathcal{D}$ zu, sodass für jedes $f : C \to C'$ in $\mathcal{C}$
\[
	\begin{tikzcd}
		F(C) \rar["F(f)"] \dar["\tau_C"]& F(C') \dar["\tau_{C'}"] \\
		G(C) \rar["G(f)"] & G(C')
	\end{tikzcd}
\]
kommutiert.
\minisec{Bemerkung}
Natürliche Transformationen lassen sich komponieren. Für $\tau : F \to G$, $\eta : G \to H$ ist $\eta \circ \tau : F \to H$ gegeben durch
\marginnote{$F,G,H : \mathcal{C} \to \mathcal{D}$}
\[
	(\eta \circ \tau)_C :=  \eta_C \circ \tau_C : F(C) \to H(C)
\]
Genauer: Für eine feste kleine Kategorie $\mathcal{C}$ und $\mathcal{D}$ bilden die Funktoren $\mathcal{C} \to \mathcal{D}$ mit den natürlichen Transformationen eine 
Kategorie $\FUN(\mathcal{C}, \mathcal{D})$.
% subsection 16 (end)

\subsection[Ausblick auf Kategorien und Funktoren in der algebraischer Topologie]{Ausblick} % (fold)
\label{sub:17}
In der algebraischen Topologie werden topologische Fragen, wie zum Beispiel, wann $\mathds{R}^n$ und $\mathds{R}^m$ homöomorph sind, in algebraische Fragen übersetzt. Eine
Möglichkeit für eine solche Übersetzung sind Funktoren von einer Kategorie von topologischen Räumen, z.B. $\TOP, \TOP_\bullet, \HTOP_\bullet$, in eine algebraisch Kategorie,
z.B. $\GRUPPEN, K\text{-}\VR, \ABELGRUPPEN, R\text{-}\MOD$. Ein Beispiel für einen solchen Funktor ist die Fundamentalgruppe $\pi_1$. Ein Nachteil der Fundamentalgruppe ist,
dass diese oft schwierig zu berechnen ist. Wir werden in dieser Vorlesung weitere Funktoren und Methoden für ihr Berechnung kennenlernen.
% subsection 17 (end)

% section 1 (end)
\newpage

\section{Kombinatorische Beschreibungen topologischer Räume} % (fold)
\label{sec:2}
\subsection[Beispiel einer kombinatorischen Beschreibung eines Raumes]{Beispiel} % (fold)
\label{sub:21}
\Index{Graphen} sind einerseits topologische Räume und andererseits kombinatorische Objekte:
\begin{align*}
	\vcenter{\hbox{\begin{tikzpicture}[scale=1.2]
		\coordinate (A) at (0,1);
		\coordinate (B) at (1,1);
		\coordinate (C) at (2,1);
		\coordinate (1) at (0,0);
		\coordinate (2) at (1,0);
		\coordinate (3) at (2,0);
		\foreach \x in {1,2,3} {
			\foreach \y in {A,B,C}{
				\draw[thick] (\x) -- (\y);
			};
		}
		\foreach \x in {A,B,C} {
			\foreach \y in {1,2,3}{
				\draw[thick] (\x) -- (\y);
			};
			\draw[fill=DodgerBlue3, thick] (\x) circle[radius=0.05] node[above]{$\x$};
		}
		\foreach \x in {1,2,3} {
			\draw[fill=DodgerBlue3,thick] (\x) circle[radius=0.05] node[below]{$\x$};
		}
	\end{tikzpicture} }}\qquad 
	&\longleftrightarrow \quad \begin{array}{l}
		V =  \set{A,B,C,1,2,3} \\
		E =  \set{ \rule{0cm}{0.3cm} { \scriptstyle \set{A,1}, \set{A,2}, \set{A,3} , \set{B,1}, \set{B,2}, \set{B,3}, \set{C,1} , \set{C,2}, \set{C,3}}}  
	\end{array}\\
	\vcenter{\hbox{\begin{tikzpicture}[scale=1.2]
		\coordinate (1) at (0,1);
		\coordinate (2) at (0.5,1.5);
		\coordinate (3) at (1,1);
		\coordinate (4) at (1,0);
		\coordinate (5) at (0,0);
		\coordinate (6) at (0.5,0.5);
		\foreach \x in {2,3,5,6} {
			\draw[thick] (1) -- (\x);
		}
		\foreach \x/\y in {2/3,3/4,3/6,4/5,4/6,5/6}{
			\draw[thick] (\x) -- (\y);
		}
		\foreach \x/\y in {1/left,2/above,3/right,4/right,5/left,6/right}{
			\draw[fill=DodgerBlue3, thick] (\x) circle[radius=0.05] node[\y]{$\x$};
		}
	\end{tikzpicture} }}\qquad
	&\longleftrightarrow \quad \begin{array}{l}
		V = \set{1,2,3,4,5,6} \\
		E = \set{ \rule{0cm}{0.3cm} {\scriptstyle \set{1,2}, \set{1,3}, \set{1,5} , \set{1,6}, \set{2,3}, \set{3,4}, \set{3,6} , \set{4,5}, \set{4,6}, \set{5,6}} }  
	\end{array}
\end{align*}
% subsection 21 (end)

\subsection[Definition: Simplizialer Komplex]{Definition} % (fold)
\label{sub:22}
Ein \Index{simplizialer Komplex} $K=(V, \Sigma)$ besteht aus einer Menge $V$ und einer Menge $\Sigma$ von nichtleeren, endlichen Teilmengen von $V$, sodass gilt
\begin{enumerate}[(i)]
	\item $\set{v} \in \Sigma$ für alle $v \in V$
	\item $\sigma \in \Sigma$, $\emptyset \not= \tau \subseteq \sigma \Longrightarrow \tau \in\Sigma.$ \hfill {\footnotesize(Abgeschlossen bzgl. Teilmengenbildung)}
\end{enumerate}
Die Elemente von $V$ heißen die \Index{Ecken} oder \Index{Vertices} von $K$. Die Elemente von $\Sigma$ heißen die \Index{Simplizes} von $K$. Enthält $\sigma \in \Sigma$ 
genau $n+1$ Elemente, so heißt $\sigma$ ein \bet{$n$-Simplex}\index{n-Simplex@$n$-Simplex}. Ist $\tau \subseteq \sigma$ mit $\sigma$ ein $n$-Simplex und $\tau$ ein $n-1$-Simplex, so heißt $\tau$ eine \Index{Seite} von $\sigma$. 

\noindent Ist $V$ geordnet, so heißt $K$ geordnet. Ist $V$ endlich, so heißt $K$ endlich.
% subsection 22 (end)

\subsection[Beispiel für simpliziale Komplexe]{Beispiel} % (fold)
\label{sub:23}
\begin{align*}
	\vcenter{\hbox{\begin{tikzpicture}[scale=1.]
		\coordinate (0) at (0,0);
		\coordinate (1) at (1,0.8);
		\coordinate (2) at (1,-0.7);
		\coordinate (3) at (2,0.2);
		\filldraw[DodgerBlue2!40] (0) -- (1) -- (2) -- cycle ;
		\foreach \x/\y in {0/1,0/2,1/2,1/3,2/3} {
			\draw[thick] (\x) -- (\y);
		};
		\foreach \x/\pos in {0/left,1/above,2/below,3/right} {
			\draw[fill=DodgerBlue3,thick] (\x) circle[radius=0.06] node[\pos]{$\x$};
		}
	\end{tikzpicture}}} \qquad 
	\begin{array}{l}
		V =  \set{0,1,2,3} \\
		E =  \set{ \rule{0cm}{0.3cm} {\emptyset, \set{0}, \set{1}, \set{2}, \set{3}, \set{0,1}, \set{0,2},\set{1,2},\set{1,3},\set{2,3},\set{0,1,2}}}  
	\end{array}
\end{align*}
Ist $\sigma$ eine endliche Menge, so heißt $\Delta^\sigma := \enbrace[\big]{\sigma, \mathcal{P}(\sigma) \setminus \set{\emptyset}}$ der $\sigma$-Simplex. Für 
$\sigma = \set{0, \ldots , n}$ schreiben wir $\Delta^n := \Delta^{\set{0, \ldots ,n}}$. Es ist
\[
	\Delta^0  = \vcenter{\hbox{\begin{tikzpicture}
		\draw[fill=DodgerBlue3,thick] (0,0) circle[radius=0.06];
	\end{tikzpicture}}} , \qquad 
	\Delta^1 = \vcenter{\hbox{\begin{tikzpicture}
		\coordinate (0) at (0,0);
		\coordinate (1) at (0.4,0.3);
		\draw[thick] (1) -- (0);
		\foreach \x in {0,1} {
			\draw[fill=DodgerBlue3,thick] (\x) circle[radius=0.06];
		}
	\end{tikzpicture}}} , \qquad 
	\Delta^2 = \vcenter{\hbox{\begin{tikzpicture}
		\coordinate (0) at (0,0);
		\coordinate (1) at (0.6,0);
		\coordinate (2) at (60:0.6);
		\draw[thick, fill=DodgerBlue2!40] (1) -- (0) -- (2) -- cycle;
		\foreach \x in {0,1,2} {
			\draw[fill=DodgerBlue3,thick] (\x) circle[radius=0.06];
		}
	\end{tikzpicture}}}, \qquad 
	\Delta^3 = 
	\vcenter{\hbox{\begin{tikzpicture}
			\coordinate (0) at (0,0);
			\coordinate (1) at (1,0);
			\coordinate (2) at (0.5,0,0.866);
			\coordinate (3) at (0.5,0.866,0.433);
			\foreach \x in {0,1,2,3}{
				\foreach \y in {0,1,2,3}{
					\foreach \z/\col in {0/DodgerBlue2,1/DodgerBlue1,2/DodgerBlue3,3/DodgerBlue2}{
					\draw[thick, fill=\col!40, fill opacity=0.4] (\x) -- (\y) -- (\z) -- cycle;
					};
				};
			};
			\foreach \x in {0,1,2,3} {
				\draw[fill=DodgerBlue3,thick] (\x) circle[radius=0.06];
			}
		\end{tikzpicture}}}
\]
% subsection 23 (end)

\subsection[Definition: Unterkomplex]{Definition} % (fold)
\label{sub:24}
Sei $K=(V, \Sigma)$ ein simplizialer Komplex. Ein \Index{Unterkomplex} von $K$ ist ein simplizialer Komplex $K_0 = (V_0, \Sigma_0)$ mit 
$V_0 \subseteq V$ und $\Sigma_0 \subseteq \Sigma$.
\minisec{Beispiel}
Sei $K=(V,\Sigma)$ ein endlicher simplizialer Komplex. Dann ist $K$ ein Unterkomplex von $\Delta^V$.
% subsection 24 (end)

\subsection[Definition: Simpliziale Abbildung]{Definition} % (fold)
\label{sub:25}
Eine \Index{simpliziale Abbildung} $f : K_1=(V_1, \Sigma_1) \to K_2(V_2, \Sigma_2)$ zwischen simplizialen Komplexen ist eine Abbildung $f : V_1 \to V_2$, sodass 
$f(\sigma_1) \in \Sigma_2$ für alle $\sigma_1 \in \Sigma_1$. 
% subsection 25 (end)

\subsection[Definition: Geometrische Realisierung]{Definition} % (fold)
\label{sub:26}
Sei $\sigma$ eine endliche Menge. Sei $\mathds{R}^\sigma = \prod_{v \in \sigma} \mathds{R}$ mit der Produkttopologie. Sei nun 
\[
	\abs*{\Delta^\sigma} := \set[x=(x_v)_{v \in \sigma} \in \mathds{R}^\sigma]{\sum_{v \in \sigma}x_v = 1 , x_v \in [0,1] \,\, \forall v \in \sigma}  
\]
Ist $\tau \subseteq \sigma$, so erhalten wir eine Abbildung $\iota_\tau^\sigma : \abs*{\Delta^\tau}  \to \abs*{\Delta^\sigma}$ indem wir 
$(x_v)_{v \in \tau} \in \abs*{\Delta^\tau}$ durch $x_v=0$ für $v \in \sigma \setminus \tau$ zu $(x_v)_{v \in \sigma}$ auffüllen. Ist $\eta \subseteq \tau \subseteq \sigma$ 
so gilt $\iota^\sigma_\eta = \iota_\tau^\sigma \circ  \iota_\eta^\tau$.

\noindent Sei $K=(V, \Sigma)$ ein simplizialer Komplex. Die \Index{geometrische Realisierung} $\abs*{K}$ von $K$ ist definiert als 
\[
	\abs*{K} := \coprod_{\sigma \in \Sigma} \set{\sigma} \times \abs*{\Delta^\sigma} \Big/ \sim
\] 
wobei $\sim$ die durch $(\tau, x) \sim \enbrace[\big]{\sigma, \iota_\tau^\sigma(x)}$ für $\tau \subseteq \sigma$ erzeugte Äquivalenzrelation ist. Versehen mit der 
Quotiententopologie ist $\abs*{K}$ ein topologischer Raum.
% subsection 26 (end)

\subsection[Beispiel: Rand des $n$-Simplizes und der Torus]{Beispiel} % (fold)
\label{sub:27}
\begin{itemize}
	\item Sei $\partial \Delta^n := \enbrace[\Big]{\set{0, \ldots ,n}, \mathcal{P}\enbrace[\big]{\set{0,\ldots ,n}} \setminus \set{\emptyset, \set{0, \ldots , n}}} = $ 
	\enquote{$\Delta^n \setminus \set{0,\ldots ,n}$}. $\partial \Delta^n$ heißt der \Index{Rand des $n$-Simplizes}. Es gilt
	\begin{align*}
		\abs*{\partial \Delta^1} &= \enspace\vcenter{\hbox{\begin{tikzpicture}
		\draw[fill] (0,0) circle[radius=0.04];
		\draw[fill] (0.15,0) circle[radius=0.04];
	\end{tikzpicture}}} \enspace\cong S^0\\ 
		\abs*{\partial \Delta^2} &=  \vcenter{\hbox{\begin{tikzpicture}
		\coordinate (0) at (0,0);
		\coordinate (1) at (0.6,0);
		\coordinate (2) at (60:0.6);
		\draw[thick] (1) -- (0) -- (2) -- cycle;
	\end{tikzpicture}}} \cong S^1 \\
		\abs*{\partial \Delta^3} &= \vcenter{\hbox{\begin{tikzpicture}
			\coordinate (0) at (0,0);
			\coordinate (1) at (1,0);
			\coordinate (2) at (0.5,0,0.866);
			\coordinate (3) at (0.5,0.866,0.433);
			\draw[sharp corners] (0) -- (1) -- (2) -- (0) -- (3) -- (2) -- (1) -- (3) -- cycle;
		\end{tikzpicture}}} \cong S^2 
	\end{align*}
	Allgemein gilt $\abs*{\partial \Delta^n} \cong S^{n-1}$.
	\item $T^2 = S^1 \times S^1 \cong 
	\vcenter{\hbox{\begin{tikzpicture}[scale=0.3]
		\draw (-3.5,0) .. controls (-3.5,2) and (-1.5,2.5) .. (0,2.5);
		\draw[xscale=-1] (-3.5,0) .. controls (-3.5,2) and (-1.5,2.5) .. (0,2.5);
		\draw[rotate=180] (-3.5,0) .. controls (-3.5,2) and (-1.5,2.5) .. (0,2.5);
		\draw[yscale=-1] (-3.5,0) .. controls (-3.5,2) and (-1.5,2.5) .. (0,2.5);

		\draw (-2,.2) .. controls (-1.5,-0.3) and (-1,-0.5) .. (0,-.5) .. controls (1,-0.5) and (1.5,-0.3) .. (2,0.2);

		\draw (-1.75,0) .. controls (-1.5,0.3) and (-1,0.5) .. (0,.5) .. controls (1,0.5) and (1.5,0.3) .. (1.75,0);
	\end{tikzpicture}}} \cong 
	\abs*{\vcenter{\hbox{\begin{tikzpicture}[>=to, scale=0.8]
		\filldraw[color=DodgerBlue2!20,] (0,0) rectangle (3,3);
		\foreach \x in {0,1,2}{
			\foreach \y in {0,1,2}{
				\draw (\x,\y) -- ++ (1,1);
			};
		};
		\draw (0,0) grid (3,3);
		\draw[very thick] (0,0) -- (0,3) -- (3,3) -- (3,0) -- cycle;
		\draw[->>, thick] (0,1) -- (0,1.5);
		\draw[->, thick] (1,3) -- (1.5,3);
		\draw[->, thick] (1,0) -- (1.5,0);
		\draw[->>, thick] (3,1) -- (3,1.5);
		\clip (-0.2,-0.2) rectangle (3.2,3.2);
	\end{tikzpicture}}}} $
\end{itemize}
% subsection 27 (end)

\subsection[Definition: Polyeder]{Definition} % (fold)
\label{sub:28}
Ein topologischer Raum $X$ heißt ein \Index{Polyeder}, falls er homöomorph zur Realisierung eines simplizialen Komplexes ist. 
% subsection 28 (end)

\subsection[Definition: Euler-Charakteristik]{Definition} % (fold)
\label{sub:29}
Sei $K=(V,\Sigma)$ ein endlicher simplizialer Komplex. Sei $a_n := \# \set[\sigma \in \Sigma]{\sigma \text{ ist ein $n$-Simplex von $K$}} $. Dann heißt 
$\chi(K) := \sum_{n \in \mathds{N}} (-1)^n a_n $ die \Index{Euler-Charakteristik} von $K$. 
\minisec{Beispiel}
\begin{align*}
	\chi \enbrace*{\vcenter{\hbox{\begin{tikzpicture}
		\coordinate (0) at (0,0);
		\coordinate (1) at (0.6,0);
		\coordinate (2) at (60:0.6);
		\draw[thick] (1) -- (0) -- (2) -- cycle;
	\end{tikzpicture}}} } &= 3 -3 = 0 \\
	\chi \enbrace*{\vcenter{\hbox{\begin{tikzpicture}
		\coordinate (0) at (0,0);
		\coordinate (1) at (0.6,0);
		\coordinate (2) at (0.6,0.6);
		\coordinate (3) at (0,0.6);
		\draw[thick] (0) -- (1) -- (2) -- (3) -- cycle;
	\end{tikzpicture}}} } &= 4-4 =0 \\
	\chi(n\text{-eck}) &= n-n = 0
\end{align*}
\todo{RevChap2}
\begin{align*}
	\chi(\partial \Delta^3 &= \text{Dreieckspyramide}) = 4 - 6 +4 = 2 \\
	\chi\enbrace*{\text{Würfel}} &= 8 -12 +6 = 2 \\
	\chi(\text{Oktaeder}) &= 6- 12  +8 = 2 
\end{align*}
% subsection 29 (end)

\subsection[Satz über die Euler-Charakteristik]{Satz} % (fold)
\label{sub:210}
Seinen $K$ und $K'$ endliche simpliziale Komplexe. Gilt $\abs*{K} \cong \abs*{K'}$, so gilt $\chi(K)= \chi(K')$.
% subsection 210 (end)
% section 2 (end)
\newpage

\section{Simpliziale Homologie} % (fold)
\label{sec:3}

\subsection[Definition: Freier $R$-Modul mit Basis $S$]{Definition} % (fold)
\label{sub:31}
Sei $S$ eine Menge.Sei $R$ ein Ring.\marginnote{Ring mit Eins} Der \bet{freie $R$-Modul mit Basis $S$}\index{freier Modul}, $R[S]$, besteht aus allen endlichen formalen $R$-Linearkombinationen
\[
	\sum_{s \in S} r_s \cdot s \quad , r_s \in R, r_s \not=0 \text{ für endliche viele }s \in S
\]
Die $R$-Modulstruktur auf $R[S]$ ist definiert durch:
\begin{align*}
	r \cdot \enbrace*{\sum_{s\in S} r_s \cdot s} &:= \sum_{s \in S} (r \cdot r_s) \cdot s \\
	\enbrace*{\sum_{s \in S} r_s \cdot s} + \enbrace*{\sum_{s \in S} r_s' s} &:= \sum_{s \in S} (r_s + r_s') \cdot s
\end{align*}
\minisec{Bemerkung}
Mittels $s= \sum_{s' \in S} \delta_{s,s'} \cdot s'$ fassen wir $S$ als Teilmenge von $R[S]$ auf.
% subsection 31 (end)

\subsection{Lemma} % (fold)
\label{sub:32}
Sei $S$ eine Menge, $R$ ein Ring.
\begin{enumerate}[1)]
	\item Ist $M$ ein $R$-Modul und $i : S \to M$ eine Abbildung, so gibt es genau eine $R$-lineare Abbildung $\varphi : R[S] \to M$, die $i$ fortsetzt.
	\item Sei $f : R[S] \to M$ $R$-linear und $p : N \twoheadrightarrow M$ $R$-linear und surjektiv. Dann gibt es $\hat{f} : R[S] \to N$ mit $p \circ \hat{f} =f$.
	\[
		\begin{tikzcd}
			\mbox{ } & N \dar["p", two heads] \\
			R[S] \rar["f"] \urar["\hat{f}", dashed] & M
		\end{tikzcd}
	\]\todo{platzsparender mit \texttt{minipage}}
\end{enumerate}
\minisec{Beweis}
\begin{enumerate}[1)]
	\item Eine solche Abbildung ist gegeben durch
	\[
		\varphi \enbrace*{\sum_{s \in S} r_s \cdot s} = \sum_{s \in S} r_s \cdot i(s) 
	\]
	Ist $\varphi'$ eine zweite, so gilt 
	\begin{align*}
		\hat{\varphi} \enbrace*{\sum_{s \in S} r_s \cdot s}  = \sum_{s \in S} r_s \cdot \hat{\varphi}(s) = \sum_{s  \in S} r_s \cdot i(s)
	\end{align*}
	\item Wähle für jedes $s \in S$ ein Urbild $j(s)\in N$ für $f(s) \in M$ unter $p : N \twoheadrightarrow M$. Nun wende 1) auf $j : S \to N$ an, um $\hat{f} : R[S] \to N$
	mit $\hat{f}(s) = j(s)$ zu erhalten. Nun ist $ p \circ \hat{f}(s) = f(s)$ für alle $s \in S$. Mit 1) folgt $p \circ \hat{f} = f$.
\end{enumerate}
% subsection 32 (end)

\subsection[Beispiel für freie und nicht-freie Moduln]{Beispiel} % (fold)
\label{sub:33}
\begin{enumerate}[(i)]
	\item Sei $K$ ein Körper, $V,M,N$ $K$-Vektorräume, $f : V \to M$, $p : N \to M$ $K$-linear, $p$ surjektiv. Dann gibt es $\hat{f} : V \to N$ mit
	$p \circ \hat{f} = f$. Wir können Lemma \ref{sub:32} benutzen, da $V$ eine Basis $B$ hat, also $V \cong K[B]$.
	\item Sei $R=\mathds{Z}$, sie $V= \nicefrac{\mathds{Z}}{2 \mathds{Z}}$, $M = \nicefrac{\mathds{Z}}{2 \mathds{Z}}$, $f= \id : V \to M$, $N= \mathds{Z}$ und 
	$p : \mathds{Z}n \to \nicefrac{\mathds{Z}}{2 \mathds{Z}}$ die Projektion mit $p(n)= n + 2 \mathds{Z}$. Dann gibt es keine $\mathds{Z}$-lineare Abbildung 
	$\hat{f} : \nicefrac{\mathds{Z}}{2 \mathds{Z}} \to \mathds{Z}$ mit $p \circ \hat{f} = f$.
	
	Ist $\hat{f} (1+ 2 \mathds{Z}) = n \in \mathds{Z}$, so folgt 
	\[
		2 n = 2 \cdot \hat{f}(1+ 2 \mathds{Z}) = \hat{f} \enbrace[\big]{2 \cdot (1+ 2 \mathds{Z})} = \hat{f} (0+ 2 \mathds{Z}) =0
	\]
	also $n=0$.
\end{enumerate}
% subsection 33 (end)

\subsection[Definition: $n$-ter Kettenmodul]{Definition} % (fold)
\label{sub:34}
Sei $K=(V, \Sigma)$ ein simplizialer Komplex. Sei $\Sigma_n := \set[\sigma \in \Sigma]{\sigma \text{ ist ein $n$-Simplex}}$. Sei $R$ ein Ring. Der 
\bet{$n$-te Kettenmodul}\index{nterKettenmodul@$n$-ter Kettenmodul} von $K$ über $R$ ist definiert als 
\[
	C_n(K;R) := R[\Sigma_n]
\]
Ist $R=\mathds{Z}$, so schreiben wir auch kurz $C_n(K) := C_n(K;\mathds{Z})$.
% subsection 34 (end)

\subsection{Definition} % (fold)
\label{sub:35}
Sei $K=(V,\Sigma)$ ein geordneter simplizialer Komplex. Sei $\sigma = \set{v_0, \ldots ,v_n}$ ein $n$-Simplex von $K$. Bezüglich der Ordnung von $K$ sei dabei 
$v_0 < v_1 < \ldots < v_n$. Sei $\tau$ ein $(n-1)$-Simplex von $K$. Definiere
\[
	\varepsilon_\sigma^\tau := \begin{cases}
		(-1)î, &\text{ falls }\tau = \set{v_0, \ldots , v_{i-1}, v_{i+1}, \ldots , v_n} \\
		0 , &\text{ sonst}
	\end{cases}
\]
\minisec{Bemerkung}
$\varepsilon^\tau_\sigma \not= 0 \iff \tau$ ist eine Seite von $\sigma$.
\minisec{Beispiel}
\[
	\text{Hier fehlt noch was } \varepsilon_{\set{0,1,2}}^{\set{0,1}} = 1 , \varepsilon_{\set{0,1,2}}^{\set{0,2}}=-1 , \varepsilon_{\set{0,1,2}}^{\set{1,2}} =1    
\]
\todo{2-Simplex mit Orientierung hinzufügen}
% subsection 35 (end)

\subsection[Definition: $n$-te Randabbildung]{Definition} % (fold)
\label{sub:36}
Sei $K=(V,\Sigma)$ ein geordneter simplizialer Komplex. Die \bet{$n$-te Randabbildung}\index{nte Randabbildung@$n$-te Randabbildung} 
\[
	\partial_n : \underset{=R[\Sigma_n]}{C_n(K;R)} \to \underset{=R[\Sigma_{n-1}]}{C_{n-1}(K;R)}
\]
ist definiert durch
\[
	\partial_n(\sigma) = \sum_{\tau \in \Sigma_{n-1}} \varepsilon_\sigma^\tau \cdot \tau \qquad \text{ für } \tau \in \Sigma_n
\]
% subsection 36 (end)

\subsection{Proposition} % (fold)
\label{sub:37}
\[
	\partial_{n-1} \circ \partial_n = 0
\]
\minisec{Beweis (mit Lemma \ref{sub:38})}
Sei $\sigma$ ein $n$-Simplex von $K$. 
\begin{align*}
	\partial_{n-1} \enbrace*{\partial_n (\sigma)} = \partial_{n-1} \enbrace*{ \sum_{\tau \in \Sigma_{n-1}} \varepsilon_\sigma^\tau \cdot \tau} = \sum_{\tau \in \Sigma_{n-1}} \varepsilon_{\sigma}^\tau \partial_{n-1} (\tau) 
	 &= \sum_{\tau \in \Sigma_{n-1}} \varepsilon_\sigma^\tau \enbrace*{ \sum_{\eta \in \Sigma_{n-2}} \varepsilon_\tau^\eta \cdot \eta} \\
	 &= \sum_{\eta \in \Sigma_{n-2}} \enbrace*{ \sum_{\tau \in \Sigma_{n-1}, \eta \subseteq \tau \subseteq \sigma} \varepsilon_\tau^\eta \varepsilon_\sigma^\tau}\eta 
	 \stackrel{\text{Lemma \ref{sub:38}}}{=} 0 \bewende
\end{align*}
% subsection 37 (end)

\subsection{Lemma} % (fold)
\label{sub:38}
Sei $\sigma$ ein $n$-Simplex. Sei $\eta \subseteq \sigma$ ein $(n-2)$-Simplex. Dann gibt es genau zwei $(n-1)$-Simplizes $\tau, \tau'$ von $K$, die eine Seite von $\sigma$
sind und $\eta$ als Seite enthalten. Es gilt
\[
	\varepsilon_\tau^\eta \varepsilon^\tau_\sigma = - \varepsilon_{\tau'}^\eta \varepsilon^{\tau'}_\sigma.
\]
\minisec{Beweis}
Sei $\sigma= \set{v_0, \ldots ,v_n}$ mit $v_0 < v_1 < \ldots < v_n$. Dann ist $\eta = \set{v_0, \ldots , v_n} \setminus \set{v_i, v_j}$ mit $i<j$. Dann sind 
$\tau = \set{v_0, \ldots , v_{i-1}, v_{i+1}, \ldots , v_n}$ und $\tau' = \set{v_0, \ldots , v_{j-1}, v_{j+1}, \ldots , v_n}$ die gesuchten $(n-1)$-Simplizes und es gilt:
\begin{align*}
	\varepsilon_\sigma^\tau = (-1)^i \enspace, \enspace \varepsilon_\sigma^{\tau'} = (-1)^j \enspace, \enspace \varepsilon^\eta_\tau = (-1)^{j-1} \enspace, \enspace 
	\varepsilon_{\tau'}^\eta = (-1)^i \bewende
\end{align*}
% subsection 38 (end)

\subsection{Notation} % (fold)
\label{sub:39}
Für $\sigma \in \Sigma_n$, $ \sigma= \set{v_0, \ldots ,v_n} $ mit $v_0 < \ldots < v_n$ schreiben wir 
\[
	\delta^j \sigma := \set{v_0 , \ldots , v_{j-1}, v_{j+1}, \ldots , v_n}  \in \Sigma_{n-1}
\]
Dann ist $\partial_n(\sigma) = \sum_{j=0}^{n} (-1)^j \delta^j \sigma$.
% subsection 39 (end)

\subsection{Bemerkung} % (fold)
\label{sub:310}
$\partial_{n-1} \circ \partial_n = 0 \iff \im \partial_n \subseteq \ker \partial_{n-1}$
% subsection 310 (end)

\subsection[Definition: $n$-ter Homologiemodul]{Definition} % (fold)
\label{sub:311}
Sei $K$ ein geordneter simplizialer Komplex. Der \bet{$n$-te Homologiemodul}\index{nter Homologiemodul@$n$-ter Homologiemodul} von $K$ über $R$ ist definiert als
\[
	H_n(K;R) := \frac{\ker \partial_n : C_n(K;R) \to C_{n-1}(K;R)}{\im \partial_{n-1} : C_{n+1}(K;R) \to C_n(K;R)} 
\]
% subsection 311 (end)
% section 3 (end)



\cleardoubleoddemptypage
\pagenumbering{Alph}
\setcounter{page}{1}
\printindex
\listoffigures
\todototoc
\listoftodos[Todo's und andere Baustellen]
\end{document}
